\subsection*{Confirmed Vulnerabilities}

\begin{table}[H]
\centering
\begin{tabular}{|l|p{6cm}|l|l|}
\hline
\textbf{ID} & \textbf{Vulnerability} & \textbf{Severity} & \textbf{Status} \\
\hline

BAC &
Broken Access Control -- Privilege Escalation via Registration &
Critical &
Confirmed \\
\hline

IDOR &
Insecure Direct Object Reference -- Unauthorized Basket Access &
Medium &
Confirmed \\
\hline

AUTH &
Weak Password Recovery &
Critical &
Confirmed \\
\hline

SQLI-01 &
SQL Injection -- Authentication Bypass &
Critical &
Confirmed \\
\hline

SQLI-02 &
SQL Injection -- Product Search and Data Extraction &
High &
Confirmed \\
\hline

XSS &
Reflected Cross-Site Scripting -- Search Parameter &
Medium &
Confirmed \\
\hline

CSRF &
Cross-Site Request Forgery -- Profile Update &
Medium &
Confirmed \\
\hline

OR &
Open Redirect -- Unvalidated Redirect Parameter &
Medium &
Confirmed \\
\hline

\end{tabular}
\caption{Summary of confirmed security vulnerabilities}
\end{table}

\subsection*{Risk Overview}

The assessment identified eight confirmed security vulnerabilities
affecting access control, authentication, application input, client-side
security, request integrity, and redirect handling.

Three findings were classified as Critical, one as High, and four as
Medium.

The Critical findings include Broken Access Control, SQL Injection --
Authentication Bypass, and Weak Password Recovery. The Broken Access
Control finding affects privilege-related actions during registration.
The Weak Password Recovery finding allows an attacker to take over an
account, including a privileged account, through a predictable
security question.

The single High-severity finding is SQL Injection -- Product Search
and Data Extraction, which allows unauthorized database interaction
without requiring authentication.

The Medium-severity findings include Reflected Cross-Site Scripting,
Insecure Direct Object Reference, Open Redirect, and Cross-Site
Request Forgery. The Reflected XSS vulnerability can allow
attacker-controlled JavaScript to execute in a victim's browser when a
malicious link is opened. The IDOR vulnerability allows access to
another user's basket by changing the basket identifier in the
request. The Open Redirect vulnerability allows an attacker-controlled
destination to redirect a victim to an external website. The CSRF
vulnerability allows an attacker-controlled webpage to submit a
profile update request while a victim is authenticated.

\subsection*{Finding Distribution}

The confirmed findings cover the following major security areas:

\begin{itemize}

    \item \textbf{Broken Access Control} -- Unauthorized privilege
    escalation through the affected registration functionality.

    \item \textbf{Insecure Direct Object Reference} -- Unauthorized
    access to another user's basket by manipulating the basket
    identifier.

    \item \textbf{Authentication} -- Weakness in the password recovery
    mechanism through the security-question process.

    \item \textbf{SQL Injection} -- Authentication bypass and
    unauthorized retrieval of application data.

    \item \textbf{Cross-Site Scripting} -- Execution of
    attacker-controlled JavaScript in a victim's browser.

    \item \textbf{Cross-Site Request Forgery} -- Unauthorized profile
    modification through a cross-site request.

    \item \textbf{Open Redirect} -- Redirect validation weakness
    allowing navigation to an external domain.

\end{itemize}

Detailed testing steps, evidence, impact, and remediation
recommendations for each confirmed finding are provided in the
following sections of this report.